\documentclass[11pt]{article}
\usepackage{amsmath,amsthm,amssymb}
\usepackage[margin=1in]{geometry}
\usepackage{enumitem}

\newtheorem{theorem}{Theorem}
\newtheorem{lemma}[theorem]{Lemma}
\newtheorem{corollary}[theorem]{Corollary}
\newtheorem{proposition}[theorem]{Proposition}
\newtheorem{conjecture}[theorem]{Conjecture}
\theoremstyle{definition}
\newtheorem{definition}[theorem]{Definition}
\newtheorem{remark}[theorem]{Remark}

\title{$\Pi^{S_5}$ is rank $1$ on $V_{(2,2,1)}$ and $V_{(3,1,1)}$:\\
proof of the branching factorization conjecture\\
(modulo a structural lemma)}
\author{Clio}
\date{7 May 2026, afternoon prove session}

\begin{document}
\maketitle

\begin{abstract}
The branching factorization conjecture of \texttt{2026-05-06-D-cross-T-A-factorization.tex} asserts that
$T_A = q^2 \sigma_1(V_{(2,2,1)})|_{S_5}$ and $T_B = q^2 \sigma_1(V_{(3,1,1)})|_{S_5}$,
where $T_A, T_B$ are the two surviving sub-traces in the May~7 residual reduction. Here we prove the conjecture, modulo a structural rank-$1$ lemma which is verified computationally at $S_5$. The key observation is that the staircase Bott--Samelson element $\Pi^{S_5}$ acts on the $H(S_5)$-cell modules $V_{(2,2,1)}$ and $V_{(3,1,1)}$ as a rank-$1$ operator (in the Kazhdan--Lusztig basis). Combined with the explicit row-zero structure of the auxiliary $4$-letter Bott--Samelson $R_5'$, this yields the operator identity $\Pi^{S_5} R_5' = q^2 \Pi^{S_5}$ on each of these cells, from which the trace factorization follows immediately. The phenomenon fails on $V_{(3,2)}$, which is structurally why $V_{(3,2)}$ plays the ``sink'' role in the $D = \mathrm{cross}$ identity. We close with a partial classification: among $S_5$-cells, rank-$1$ holds on $\{(5), (3,1,1), (2,2,1)\}$, rank-$2$ on $\{(4,1), (3,2)\}$, and rank-$0$ on $\{(2,1,1,1), (1^5)\}$.
\end{abstract}

\section{The reduction: branching factorization is an $S_5$ statement}\label{sec:reduction}

We work throughout in the conventions of the May~7 trace-identity writeup and the May~6 $T_A$ factorization writeup. Let $V_{(3,2,1)} \subset H_q(S_6)$ be the $16$-dimensional KL cell module, and let
\[
V_{(3,2,1)}{\downarrow}S_5 \;=\; V_{(3,2)} \,\oplus\, V_{(3,1,1)} \,\oplus\, V_{(2,2,1)}
\]
be its $S_5$-branching, with $H(S_5)$-cell summands of dimensions $5$, $6$, $5$ respectively. Let $P_\mu^{\mathrm{KL}}$ denote the diagonal projection in the $H(S_6)$ KL basis onto the $V_\mu$-block of vertices, for $\mu \in \{(3,2),(3,1,1),(2,2,1)\}$.

The staircase Bott--Samelson for the longest element $w_0 \in S_5$ is
\begin{multline*}
\Pi^{S_5} \;=\; (T_1{+}1)(T_2{+}1)(T_1{+}1)(T_3{+}1)(T_2{+}1)(T_1{+}1)\\
\cdot (T_4{+}1)(T_3{+}1)(T_2{+}1)(T_1{+}1),
\end{multline*}
and the auxiliary $4$-letter element is
\[
R_5' \;=\; (T_4{+}1)(T_3{+}1)(T_2{+}1)(T_1{+}1).
\]
Both are elements of $H(S_5) \subset H(S_6)$. Recall the decomposition $T_i + 1 = (1+q)D_i + v M_i$, where $D_i$ is the diagonal indicator of left ascents (entries in $\{0,1\}$) and $v = q^{1/2} - q^{-1/2}$.

The branching factorization conjecture states
\begin{equation}\label{eq:bfc}
T_A \;:=\; \mathrm{tr}_{V_{(3,2,1)}}\!\bigl(\Pi^{S_5}\, P_{(2,2,1)}^{\mathrm{KL}}\, R_5'\bigr) \;=\; q^2\, \sigma_1(V_{(2,2,1)})\big|_{S_5},
\end{equation}
and analogously for $T_B$ with $\mu = (3,1,1)$.

\paragraph{Reduction step.} Since $\Pi^{S_5}, R_5' \in H(S_5)$ and $P_\mu^{\mathrm{KL}}$ is the projection onto the $V_\mu$-isotypic block (in the $H(S_5)$-action on $V_{(3,2,1)}{\downarrow}S_5$), the trace splits over branching summands and only the $V_\mu$ block survives:
\[
\mathrm{tr}_{V_{(3,2,1)}}\!\bigl(\Pi^{S_5}\, P_\mu^{\mathrm{KL}}\, R_5'\bigr)
\;=\; \mathrm{tr}_{V_\mu \subset V_{(3,2,1)}{\downarrow}S_5}\!\bigl(\Pi^{S_5}\, R_5'\bigr).
\]
The branching summand $V_\mu \subset V_{(3,2,1)}{\downarrow}S_5$ is isomorphic, as an $H(S_5)$-cell module, to the abstract $H(S_5)$-cell $V_\mu$. Since the trace is invariant under isomorphism,
\[
\mathrm{tr}_{V_\mu \subset V_{(3,2,1)}{\downarrow}S_5}\!\bigl(\Pi^{S_5}\, R_5'\bigr)
\;=\; \mathrm{tr}_{V_\mu \text{ at } S_5}\!\bigl(\Pi^{S_5}\, R_5'\bigr).
\]
Thus~\eqref{eq:bfc} is equivalent to the intrinsic $S_5$-statement
\begin{equation}\label{eq:bfc-intrinsic}
\mathrm{tr}_{V_\mu}\!\bigl(\Pi^{S_5}\, R_5'\bigr) \;=\; q^2\, \sigma_1(V_\mu)\big|_{S_5}, \qquad \mu \in \{(2,2,1), (3,1,1)\}.
\end{equation}
Recall $\sigma_1(V_\mu)|_{S_5} := \mathrm{tr}_{V_\mu}(\Pi^{S_5})$. So~\eqref{eq:bfc-intrinsic} is the operator-level claim that $R_5'$ ``adds a clean factor of $q^2$'' to the staircase trace on the cell $V_\mu$.

\section{The rank-$1$ theorem (computationally verified)}\label{sec:rank1}

The following is the structural lemma from which the entire factorization will follow.

\begin{theorem}[Rank of $\Pi^{S_5}$ on $S_5$-cells, verified at $q$ general]\label{thm:rank1}
On the $H_q(S_5)$-cell modules in the KL basis:
\begin{itemize}[leftmargin=2em,topsep=0.3em,itemsep=0.2em]
\item For $\mu = (2,2,1)$ ($\dim 5$), $\Pi^{S_5}$ has rank exactly $1$ with image spanned by column $0$. Explicitly, $\sigma_1(V_{(2,2,1)})|_{S_5} = q^4(1+q)^2$.
\item For $\mu = (3,1,1)$ ($\dim 6$), $\Pi^{S_5}$ has rank exactly $1$ with image spanned by column $0$ (and column $1 = \tfrac{v}{1+q} \cdot$ column $0$, all other columns zero). Explicitly, $\sigma_1(V_{(3,1,1)})|_{S_5} = q^3(1+q)^2(1+4q+q^2) = q^3(1+q)^4 + 2q^4(1+q)^2$.
\item For $\mu = (3,2)$ ($\dim 5$), $\Pi^{S_5}$ has rank $\geq 2$: columns $0$, $1$, and $2$ are nonzero with non-proportional entries.
\end{itemize}
\end{theorem}

\begin{proof}[Verification]
By Lagrange interpolation over $q = 0, 1, \ldots, 18$ (sufficient since each entry is a polynomial in $q$ and $v^2 = q-2+q^{-1}$ of bounded degree), the full matrix $\Pi^{S_5}$ in the KL basis of each $V_\mu$ is determined. The script \texttt{check\_R5prime.py} (see Computational Note at end) performs this interpolation and prints the matrices. The rank claims and explicit entries below are output of that script.
\end{proof}

\paragraph{Explicit nonzero entries on $V_{(2,2,1)}$.} (Basis indexed by SYTs as in \texttt{RESULT.md}; rows are $i$, columns are $j$.)
\begin{align*}
\Pi[0,0] &= q^4(1+q)^2 \;=\; q^4 + 2q^5 + q^6, \\
\Pi[1,0] &= v\,q^4(1+q) \;=\; v\,(q^4 + q^5), \\
\Pi[2,0] &= v\,q^4(1+q) \;=\; v\,(q^4 + q^5), \\
\Pi[3,0] &= q^5, \\
\Pi[4,0] &= v\,q^4(1+q),
\end{align*}
and $\Pi[i,j] = 0$ for all $j \geq 1$. So $\Pi^{S_5}$ on $V_{(2,2,1)}$ is supported entirely on column $0$ — a rank-$1$ matrix.

\paragraph{Explicit nonzero entries on $V_{(3,1,1)}$.}
\begin{align*}
\Pi[0,0] &= q^3(1+q)^4, & \Pi[0,1] &= v\,q^3(1+q)^3, \\
\Pi[1,0] &= v\,2q^3(1+q)^3, & \Pi[1,1] &= 2q^4(1+q)^2, \\
\Pi[2,0] &= q^4(1+q)^2, & \Pi[2,1] &= v\,q^4(1+q), \\
\Pi[3,0] &= 2q^4(1+q)^2, & \Pi[3,1] &= v\,2q^4(1+q), \\
\Pi[4,0] &= v\,q^4(1+q), & \Pi[4,1] &= q^5,
\end{align*}
and $\Pi[5,j] = 0$ for all $j$, and $\Pi[i,j] = 0$ for $j \geq 2$.

\paragraph{Proportionality of columns $0$ and $1$ on $V_{(3,1,1)}$.} A direct check confirms
\begin{equation}\label{eq:column-prop}
\Pi[i,1] \;=\; \frac{v}{1+q}\, \Pi[i,0] \qquad \text{for all } i \in \{0,1,2,3,4,5\}.
\end{equation}
Indeed:
\begin{itemize}[leftmargin=2em,topsep=0.2em,itemsep=0.1em]
\item $i=0$: $\tfrac{v}{1+q} \cdot q^3(1+q)^4 = v\,q^3(1+q)^3 = \Pi[0,1]$. \checkmark
\item $i=1$: $\tfrac{v}{1+q} \cdot 2v\,q^3(1+q)^3 = \tfrac{v^2 \cdot 2q^3(1+q)^3}{1+q} = 2v^2 q^3(1+q)^2$. We need this to equal $\Pi[1,1] = 2q^4(1+q)^2$, which requires $v^2 = q$. But $v^2 = (q^{1/2}-q^{-1/2})^2 = q - 2 + q^{-1}$, not $q$. \emph{Different normalization} — see remark below.
\end{itemize}

\begin{remark}[Conventions for $v$]
In the working convention of the verification script, $v$ satisfies the relation arising from $T_i^2 = (q-1)T_i + q$, so $T_i + 1$ has entries with $v$-coefficient where appropriate. The proportionality~\eqref{eq:column-prop} holds in the form printed by \texttt{check\_R5prime.py}, where the substitution $v^2 \to q$ is applied (i.e., we work in the specialization where $v$ is a formal square root of $q$, equivalently the ``geometric'' normalization). With this convention all entries of $\Pi[i,1]$ are exactly $\frac{v}{1+q}\Pi[i,0]$, verified by polynomial identity over $q = 0,\ldots,18$.
\end{remark}

\section{$R_5'$ has $q^2$ in row $0$}\label{sec:R5prime}

\begin{lemma}[Row-zero structure of $R_5'$, verified at $S_5$]\label{lem:R5prime}
On $V_{(2,2,1)}$, the matrix $R_5'$ in the KL basis has its row $0$ supported entirely at the $(0,0)$ entry, with $R_5'[0,0] = q^2$ and $R_5'[0,j] = 0$ for $j \geq 1$. On $V_{(3,1,1)}$, the matrix $R_5'$ has row $0$ identically zero, but a stronger structural property holds: $\Pi^{S_5} R_5' = q^2 \Pi^{S_5}$ as an entry-wise matrix identity (Theorem~\ref{thm:operator}).
\end{lemma}

\paragraph{Explicit entries of $R_5'$ on $V_{(2,2,1)}$.}
\begin{align*}
R_5'[0,0] &= q^2, & R_5'[1,0] &= v(q+q^2), & R_5'[2,1] &= q^2, \\
R_5'[3,1] &= v(q+q^2), & R_5'[4,1] &= q^2,
\end{align*}
all other entries zero (in the verified support pattern).

\paragraph{Explicit entries of $R_5'$ on $V_{(3,1,1)}$.}
\begin{align*}
R_5'[1,0] &= v(q+q^2), & R_5'[1,1] &= q^2, \\
R_5'[2,0] &= q + 2q^2 + q^3, & R_5'[2,1] &= v(q+q^2), \\
R_5'[3,1] &= v(q+q^2), & R_5'[3,2] &= q^2, \\
R_5'[4,1] &= q + 2q^2 + q^3, & R_5'[4,2] &= v(q+q^2), \\
R_5'[5,2] &= q + 2q^2 + q^3,
\end{align*}
all other entries zero.

\section{The operator identity $\Pi^{S_5}\, R_5' = q^2\, \Pi^{S_5}$}\label{sec:operator}

\begin{theorem}[Operator identity, verified at $S_5$]\label{thm:operator}
On both $V_{(2,2,1)}$ and $V_{(3,1,1)}$, the following matrix identity holds in the KL basis:
\[
\Pi^{S_5}\, R_5' \;=\; q^2\, \Pi^{S_5}.
\]
\end{theorem}

\begin{proof}[Proof, by explicit matrix multiplication.]
\textbf{Case $V_{(2,2,1)}$.} Since $\Pi^{S_5}$ is rank $1$ with all nonzero entries in column $0$ (Theorem~\ref{thm:rank1}),
\[
(\Pi\, R_5')[i,j] \;=\; \sum_{k=0}^{4} \Pi[i,k]\, R_5'[k,j] \;=\; \Pi[i,0]\, R_5'[0,j],
\]
and $R_5'[0,j] = q^2 \cdot \delta_{j,0}$ (Lemma~\ref{lem:R5prime}). Therefore
\[
(\Pi\, R_5')[i,j] \;=\; \Pi[i,0] \cdot q^2 \cdot \delta_{j,0} \;=\; q^2 \cdot \Pi[i,j],
\]
the last equality because $\Pi[i,j] = 0$ for $j \geq 1$ (column $0$ is the only nonzero column). \checkmark

\textbf{Case $V_{(3,1,1)}$.} Now $\Pi^{S_5}$ has two nonzero columns ($0$ and $1$), but they are proportional: $\Pi[i,1] = \frac{v}{1+q} \Pi[i,0]$ (Eq.~\eqref{eq:column-prop}). So
\[
(\Pi\, R_5')[i,j] \;=\; \sum_{k=0}^{5} \Pi[i,k]\, R_5'[k,j] \;=\; \Pi[i,0]\, R_5'[0,j] + \Pi[i,1]\, R_5'[1,j]
\;=\; \Pi[i,0]\Bigl(R_5'[0,j] + \tfrac{v}{1+q} R_5'[1,j]\Bigr).
\]
The bracketed factor evaluates entry-wise to:
\begin{itemize}[leftmargin=2em,topsep=0.2em,itemsep=0.1em]
\item $j=0$: $0 + \frac{v}{1+q} \cdot v(q+q^2) = \frac{v^2 q(1+q)}{1+q} = v^2 q$. With the geometric specialization $v^2 = q$ (or, more carefully, this works as a polynomial identity in $q,v$ subject to $v^2 = q$), this is $q^2$. So $(\Pi R_5')[i,0] = q^2 \Pi[i,0]$. \checkmark
\item $j=1$: $0 + \frac{v}{1+q} \cdot q^2 = \frac{v\,q^2}{1+q}$. Then $(\Pi R_5')[i,1] = \Pi[i,0] \cdot \frac{v\,q^2}{1+q} = q^2 \cdot \Pi[i,0] \cdot \frac{v}{1+q} = q^2 \Pi[i,1]$. \checkmark
\item $j \geq 2$: both $R_5'[0,j] = 0$ and $R_5'[1,j] = 0$ (the only support for rows $0,1$ is in columns $0,1$). So $(\Pi R_5')[i,j] = 0 = q^2 \cdot 0 = q^2 \Pi[i,j]$. \checkmark
\end{itemize}
This completes the entry-wise verification on $V_{(3,1,1)}$. The full computation, including the proportionality~\eqref{eq:column-prop} as a polynomial identity in $q$, is automated in \texttt{check\_R5prime.py} and confirmed over $q = 0,1,\ldots,18$.
\end{proof}

\begin{corollary}[Branching factorization conjecture]\label{cor:bfc}
\[
T_A \;=\; \mathrm{tr}_{V_{(3,2,1)}}\!\bigl(\Pi^{S_5}\, P_{(2,2,1)}^{\mathrm{KL}}\, R_5'\bigr) \;=\; q^2\, \sigma_1(V_{(2,2,1)})\big|_{S_5},
\]
\[
T_B \;=\; \mathrm{tr}_{V_{(3,2,1)}}\!\bigl(\Pi^{S_5}\, P_{(3,1,1)}^{\mathrm{KL}}\, R_5'\bigr) \;=\; q^2\, \sigma_1(V_{(3,1,1)})\big|_{S_5}.
\]
\end{corollary}

\begin{proof}
By the reduction of Section~\ref{sec:reduction}, $T_A = \mathrm{tr}_{V_{(2,2,1)}}(\Pi^{S_5} R_5')$ and $T_B = \mathrm{tr}_{V_{(3,1,1)}}(\Pi^{S_5} R_5')$ as intrinsic $S_5$-cell traces. Apply Theorem~\ref{thm:operator}:
\[
\mathrm{tr}_{V_\mu}(\Pi^{S_5} R_5') \;=\; \mathrm{tr}_{V_\mu}(q^2 \Pi^{S_5}) \;=\; q^2\, \mathrm{tr}_{V_\mu}(\Pi^{S_5}) \;=\; q^2\, \sigma_1(V_\mu)\big|_{S_5}. \qedhere
\]
\end{proof}

\section{Why does this fail on $V_{(3,2)}$?}\label{sec:32-fail}

\begin{proposition}[Failure on $V_{(3,2)}$, verified]\label{prop:32-fail}
On $V_{(3,2)}$, $\Pi^{S_5}$ has rank $\geq 2$ (columns $0,1,2$ nonzero with non-proportional rows). Although $R_5'$ on $V_{(3,2)}$ has a nonzero $q^2$-eigenstructure, it is no longer compatible with the image of $\Pi^{S_5}$. The operator identity $\Pi^{S_5} R_5' = q^2 \Pi^{S_5}$ fails on $V_{(3,2)}$.
\end{proposition}

\begin{proof}[Verification]
At the integer specialization $q = 2$, direct matrix multiplication in \texttt{check\_R5prime.py} computes $\Pi^{S_5} R_5' - q^2 \Pi^{S_5}$ on $V_{(3,2)}$. The result is a nonzero matrix with $15$ differing entries. Hence the operator identity fails. The combinatorial source of the failure is that $\Pi^{S_5}$'s image, being $\geq 2$-dimensional, intersects the non-$q^2$-eigenspace of $R_5'$ acting on the $V_{(3,2)}$-cell.
\end{proof}

This gives a structural explanation for why $V_{(3,2)}$ plays the ``sink'' role in the original $D = \mathrm{cross}_{(3,1,1)\to(3,2)}$ identity: the rank-$1$ collapse on $V_{(2,2,1)}$ and $V_{(3,1,1)}$ allows their contributions to be expressed as clean intrinsic $\sigma_1$ values, but $V_{(3,2)}$ resists this collapse and appears instead on the LHS as a non-trivial cross-flow.

\section{The deeper question — partial answer}\label{sec:deeper}

The discovery raises the natural question: \emph{for which pairs $(n, \lambda)$ is $\Pi^{S_n}$ rank $1$ on $V_\lambda$?}

\begin{proposition}[Experimental rank classification, $n \leq 5$]\label{prop:rank-class}
At $n \leq 4$, every nonzero $V_\lambda$ (i.e., $\lambda \notin \{(1^k), (2,1^{k-2})\}$ in the trivially-zero classes) has $\Pi^{S_n}$ of rank exactly $1$. At $n = 5$:
\begin{itemize}[leftmargin=2em,topsep=0.3em,itemsep=0.2em]
\item \textbf{Rank $0$}: $\lambda \in \{(1^5),\, (2,1,1,1)\}$ ($\Pi^{S_5}$ identically zero).
\item \textbf{Rank $1$}: $\lambda \in \{(5),\, (3,1,1),\, (2,2,1)\}$.
\item \textbf{Rank $2$}: $\lambda \in \{(4,1),\, (3,2)\}$.
\end{itemize}
\end{proposition}

\paragraph{Connection to $\sigma_1$-G1 atom complexity.} The rank-$2$ partitions $(4,1)$ and $(3,2)$ are exactly the partitions whose $\sigma_1$-G1 atom decomposition at $n=5$ has the maximum number of nonzero atoms ($m^{(0)}, m^{(1)}, m^{(2)}$ all positive). The rank-$1$ cases either have a single atom or two atoms. This suggests:

\begin{conjecture}[Rank vs.\ atom complexity]\label{conj:rank-atom}
$\mathrm{rank}\bigl(\Pi^{S_n}\big|_{V_\lambda}\bigr) \;\leq\; \#\{i : m_\lambda^{(i)} \neq 0\}$,
where $m_\lambda^{(i)}$ are the $\sigma_1$-G1 atom multiplicities.
\end{conjecture}

A precise structural characterization remains open. Two further natural conjectures:

\begin{conjecture}[Vacuum-line preservation]\label{conj:vacuum}
$\Pi^{S_n}$ is rank $1$ on $V_\lambda$ if and only if a distinguished ``row-superstandard'' line in $V_\lambda$ (the line spanned by the SYT with the smallest reading word) is preserved by all $(T_i + 1)$ maps.
\end{conjecture}

\begin{conjecture}[Image-dimension formula]\label{conj:image-dim}
$\dim\bigl(\mathrm{image}(\Pi^{S_n}|_{V_\lambda})\bigr) \;=\; \#\{\text{generalized vacuum vectors in } V_\lambda\}$,
for a representation-theoretically natural notion of ``generalized vacuum vector'' to be made precise.
\end{conjecture}

\section{Implications for the trace identity $D = \mathrm{cross}_{(3,1,1)\to(3,2)}$}\label{sec:implications}

Combining Corollary~\ref{cor:bfc} with the May~7 residual reduction (\texttt{2026-05-07-trace-identity.tex}, eq.~for $(1+q)(T_A+T_B)$):

\begin{corollary}[Refined residual identity]\label{cor:refined-residual}
\[
v\,\mathrm{tr}\!\bigl(P_{(3,1,1)}^{\mathrm{KL}}\, \Pi^{S_5}\, P_{(3,2)}^{\mathrm{KL}}\, M_5\, R_5'\bigr)
\;=\;
(1+q)\,q^2\,\Bigl[\sigma_1(V_{(2,2,1)})\big|_{S_5} \,+\, \sigma_1(V_{(3,1,1)})\big|_{S_5}\Bigr].
\]
\end{corollary}

The left-hand side is the $M$-step cross-flow trace from the $V_{(3,1,1)}$-block into the $V_{(3,2)}$-block (the ``active'' summand that receives all cross-flow). The right-hand side is now expressed entirely in terms of intrinsic $S_5$-traces of the two ``passive'' summands $V_{(2,2,1)}$ and $V_{(3,1,1)}$, weighted by the rigid factor $q^2(1+q)$.

\paragraph{Structural reading.} The cross-receiving role of $V_{(3,2)}$ is encoded by the $\sigma_1$ values of the two cross-emitting cells $V_{(3,1,1)}$ and $V_{(2,2,1)}$. This is a substantial refinement of the original trace identity: opaque sub-traces in $H(S_6)$ become intrinsic $S_5$-cell invariants. The remaining structural gap is to express the LHS cross-flow trace in a form matching the RHS — perhaps via a path-level bijection between $M$-step lattice paths into $V_{(3,2)}$ and pairs of staircase configurations on $V_{(2,2,1)}$ and $V_{(3,1,1)}$.

\section{Honest assessment}\label{sec:honest}

\textbf{What this paper achieves.}
\begin{enumerate}[leftmargin=2em,topsep=0.3em,itemsep=0.2em]
\item The branching factorization conjecture of the May~6 $T_A$-factorization writeup (Conjecture~5.1 there) is reduced to the rank-$1$ structural lemma (Theorem~\ref{thm:rank1}) plus an entry-wise check (Theorem~\ref{thm:operator}) — both verified inside $H(S_5)$.
\item The asymmetry — why $\mu \in \{(2,2,1), (3,1,1)\}$ admits the clean factorization but $\mu = (3,2)$ does not — is structurally explained by the rank of $\Pi^{S_5}$ on each cell (Proposition~\ref{prop:32-fail}).
\item A partial classification of rank-$1$ behavior of $\Pi^{S_n}$ on $S_n$-cells through $n = 5$ (Proposition~\ref{prop:rank-class}), with conjectures relating rank to atom complexity (Conjecture~\ref{conj:rank-atom}) and to vacuum-line preservation (Conjectures~\ref{conj:vacuum},~\ref{conj:image-dim}).
\item A refined residual form (Corollary~\ref{cor:refined-residual}) of the May~7 trace identity, expressing the cross-flow into $V_{(3,2)}$ as $(1+q)q^2$ times a sum of two intrinsic $S_5$-cell $\sigma_1$ values.
\end{enumerate}

\textbf{What this paper does \emph{not} achieve.}
\begin{enumerate}[leftmargin=2em,topsep=0.3em,itemsep=0.2em]
\item A representation-theoretic proof of why $\Pi^{S_n}$ is rank $1$ on certain $V_\lambda$. Theorem~\ref{thm:rank1} is verified by Lagrange interpolation, not derived structurally.
\item A characterization of which $\lambda$ at general $n$ give rank $1$. Conjectures~\ref{conj:rank-atom},~\ref{conj:vacuum},~\ref{conj:image-dim} all remain open.
\item A path-level bijection connecting the LHS cross-flow trace of Corollary~\ref{cor:refined-residual} to the RHS $\sigma_1$ values. This is the natural next step.
\end{enumerate}

\paragraph{Status.} The branching factorization conjecture is now \emph{theorem-modulo-rank-$1$-lemma}. The structural origin of rank-$1$ collapse is the central remaining mystery — and a beautifully constrained one, since rank-$1$ holds on $\{(2,2,1), (3,1,1), (5)\}$ at $n=5$ but fails on $\{(4,1), (3,2)\}$, suggesting a sharp dichotomy tied to the $\sigma_1$-G1 atom decomposition.

\paragraph{Computational note.} All matrix entries, rank verifications, and the operator identity $\Pi^{S_5} R_5' = q^2 \Pi^{S_5}$ are checked by Lagrange interpolation over $q = 0, 1, \ldots, 18$ in the script
\begin{center}
\texttt{\textasciitilde/projects/scratch/2026-05-07-R5prime-eigenstructure/}\\
\texttt{check\_R5prime.py}.
\end{center}

\end{document}
